\documentclass[11pt,a4paper]{article}
\usepackage[utf8]{inputenc}
\usepackage[T1]{fontenc}
\usepackage[spanish,es-minimal,es-noshorthands]{babel}
\AtBeginDocument{\renewcommand{\sin}{\operatorname{sen}}\renewcommand{\arcsin}{\operatorname{arcsen}}}
\usepackage{lmodern}
\usepackage[sfdefault,lining]{FiraSans}
\usepackage[a4paper,top=2.3cm,bottom=2.2cm,left=1.8cm,right=1.8cm,headheight=1.3cm,headsep=0.4cm,footskip=1cm]{geometry}
\usepackage{amsmath,amssymb}
\usepackage{icomma}
\usepackage{graphicx}
\usepackage[table]{xcolor}
\usepackage{tikz}
\usepackage{fontawesome5}
\usepackage{ccicons}
\usepackage{fancyhdr}
\usepackage{lastpage}
\usepackage{tasks}
\usepackage{needspace}
\usepackage[most]{tcolorbox}
\usepackage{array}
\usepackage[hidelinks]{hyperref}
\AtBeginDocument{\hypersetup{pdftitle={Ficha de ejercicios de la unidad 5: Cálculo de derivadas. Derivabilidad de una función. 2º Bach.},pdfsubject={Matemáticas CCSS 2º Bach.},pdfkeywords={derivadas, tasa de variación media, definición de derivada, recta tangente, reglas de derivación, regla de la cadena, derivadas laterales, derivabilidad, funciones a trozos, parámetros, PAU}}}

\definecolor{navy}{HTML}{1B2748}
\definecolor{accent}{HTML}{E8603A}
\definecolor{teal}{HTML}{1F9E90}
\definecolor{soft}{HTML}{EEF1F8}
\definecolor{muted}{HTML}{5D6685}

\setlength{\parindent}{0pt}
\setlength{\parskip}{3pt}
\settasks{label=\alph*), label-format=\bfseries\color{navy}, label-width=1.3em, item-indent=1.8em, column-sep=0.6em, after-item-skip=5pt, before-skip=3pt, after-skip=4pt}

% ---------- capçalera i peu ----------
\newcommand{\llicencia}{\href{https://creativecommons.org/licenses/by-nc-sa/4.0/deed.es}{\ccbyncsa}\ \ Licencia Creative Commons Atribución-NoComercial-CompartirIgual 4.0 Internacional (CC BY-NC-SA 4.0)}
\pagestyle{fancy}
\fancyhf{}
\fancyhead[L]{\raisebox{-0.35\height}{\includegraphics[height=0.95cm]{logo-gva.png}}}
\fancyhead[C]{\small\color{muted}Ficha de ejercicios · Unidad 5: Cálculo de derivadas · 2\textsuperscript{o} Bach.}
\fancyhead[R]{\raisebox{-0.35\height}{\includegraphics[height=0.95cm]{logo-vermellar.png}}}
\fancyfoot[L]{\scriptsize\color{muted}\llicencia}
\fancyfoot[R]{\small\color{muted}\thepage}
\renewcommand{\headrulewidth}{0.4pt}
\renewcommand{\headrule}{\color{muted!40}\hrule height \headrulewidth}
\fancypagestyle{primera}{\fancyhf{}\renewcommand{\headrulewidth}{0pt}
  \fancyfoot[L]{\scriptsize\color{muted}\llicencia}
  \fancyfoot[R]{\small\color{muted}\thepage}}

% ---------- nivells i exercicis ----------
\newcommand{\estrelles}[1]{{\color{accent}\ifcase#1\or\faStar\faStar[regular]\faStar[regular]\or\faStar\faStar\faStar[regular]\or\faStar\faStar\faStar\fi}}
\newcounter{exer}
\newcommand{\ex}[1]{\par\needspace{4\baselineskip}\medskip\refstepcounter{exer}\noindent{\color{navy}\bfseries\theexer.}\ \ }

\newcommand{\seccio}[2]{\par\needspace{10\baselineskip}\bigskip
  \begin{tcolorbox}[colback=navy,colframe=navy,coltext=white,arc=2mm,boxrule=0pt,left=3mm,right=3mm,top=1.2mm,bottom=1.2mm]
  {\large\bfseries #1\quad #2}
  \end{tcolorbox}\vspace{-2pt}}


% quadre de recordatori (només a la fitxa d'exercicis)
\newcommand{\recorda}[1]{\par\smallskip\begin{tcolorbox}[colback=accent!7,colframe=accent!55,arc=2mm,boxrule=0.5pt,left=2.5mm,right=2.5mm,top=1mm,bottom=1mm]\small{\color{accent}\faLightbulb[regular]}\ \textbf{Recuerda:}\ #1\end{tcolorbox}}
% etiqueta dels exercicis de la PAU d'altres comunitats
\newcommand{\pau}[1]{{\color{teal}\small\bfseries\faUniversity\ PAU · #1.}\ }
\newcommand{\res}[1]{#1}\newcommand{\rest}[1]{#1}

% =====================================================================
% Capçalera FentMates de les fitxes: un color per a cada curs.
% Els colors són els de la pàgina inicial (LEVELS a assets/js/materials.js);
% el logo ve de _eines/logo/pdf/ (python3 _eines/logo/genera.py).
%
% Ús, al final del preàmbul de la unitat:
%   % =====================================================================
% Capçalera FentMates de les fitxes: un color per a cada curs.
% Els colors són els de la pàgina inicial (LEVELS a assets/js/materials.js);
% el logo ve de _eines/logo/pdf/ (python3 _eines/logo/genera.py).
%
% Ús, al final del preàmbul de la unitat:
%   % =====================================================================
% Capçalera FentMates de les fitxes: un color per a cada curs.
% Els colors són els de la pàgina inicial (LEVELS a assets/js/materials.js);
% el logo ve de _eines/logo/pdf/ (python3 _eines/logo/genera.py).
%
% Ús, al final del preàmbul de la unitat:
%   \input{../comu/fentmates.tex}
%   \fmcurs{1eso}
% i, al començament del document:
%   \capcalera{Matemàtiques}{1r ESO}{Fitxa d'exercicis · Unitat 2}{Divisibilitat}
% Opcional abans de \input: \def\fmneed{14} (espai mínim abans de cada secció).
% =====================================================================
\makeatletter
\def\fm@logodir{../../logo/pdf}
\providecommand{\fmneed}{10}
\def\fm@cursnom{}

% clau del curs, color i color del text damunt del color
\newcommand{\fm@curs}[3]{\@namedef{fm@color@#1}{#2}\@namedef{fm@text@#1}{#3}}
\fm@curs{1eso}{E6B84A}{navy}
\fm@curs{2eso}{3B82F6}{white}
\fm@curs{3eso}{FF7A50}{white}
\fm@curs{4eso-a}{A78BFA}{white}
\fm@curs{4eso-b}{6D28D9}{white}
\fm@curs{1batx-mi}{9B1C1C}{white}
\fm@curs{1batx-ccss}{EF6B6B}{white}
\fm@curs{2batx-mii}{1E6B47}{white}
\fm@curs{2batx-ccss}{2BB3A3}{white}

\newcommand{\fmcurs}[1]{%
  \@ifundefined{fm@color@#1}{\PackageError{fentmates}{Curs desconegut: #1}{}}{}%
  \edef\fm@tmp{\noexpand\definecolor{curs}{HTML}{\@nameuse{fm@color@#1}}%
               \noexpand\colorlet{curstext}{\@nameuse{fm@text@#1}}}\fm@tmp
  \colorlet{cursfosc}{curs!72!black}%
  \def\fm@logo{\fm@logodir/vertical-#1.pdf}%
  \def\fm@marca{\fm@logodir/marca-#1.pdf}}

% ---------- capçalera de la primera pàgina ----------
% #1 matèria (petit), #2 curs, #3 tipus de document i unitat, #4 títol
\newcommand{\capcalera}[4]{%
  \gdef\fm@cursnom{#2}%
  \par\noindent
  \begin{tcolorbox}[colback=curs,colframe=curs,boxrule=0pt,arc=2.2mm,width=3.5cm,height=2.15cm,
      valign=center,left=3mm,right=2mm,top=1mm,bottom=1mm,nobeforeafter,box align=center]
    \color{curstext}{\footnotesize #1}\par\vspace{1pt}{\fontsize{19}{21}\selectfont\bfseries #2\par}%
  \end{tcolorbox}\hspace{5mm}%
  \begin{minipage}[c]{\dimexpr\linewidth-3.5cm-5mm-2.9cm\relax}
    {\color{cursfosc}\small\bfseries\MakeUppercase{#3}}\par\vspace{2pt}
    {\color{navy}\fontsize{20}{23}\selectfont\bfseries #4\par}
  \end{minipage}\hfill
  \begin{minipage}[c]{2.6cm}\raggedleft\includegraphics[width=2.6cm]{\fm@logo}\end{minipage}\par
  \vspace{2pt}}

% ---------- barres de secció del color del curs ----------
\renewcommand{\seccio}[2]{\par\needspace{\fmneed\baselineskip}\bigskip
  \begin{tcolorbox}[colback=curs,colframe=curs,coltext=curstext,arc=2mm,boxrule=0pt,left=3mm,right=3mm,top=1.2mm,bottom=1.2mm]
  {\large\bfseries #1\quad #2}
  \end{tcolorbox}\vspace{-2pt}}

% ---------- estrelles de dificultat del color del curs (a Batxillerat no se'n posen) ----------
\@ifundefined{estrelles}{}{%
  \renewcommand{\estrelles}[1]{{\color{curs}\ifcase#1\or\faStar\faStar[regular]\faStar[regular]\or\faStar\faStar\faStar[regular]\or\faStar\faStar\faStar\fi}}}

% ---------- capçalera i peu de les pàgines ----------
% primera pàgina: logos institucionals al peu
\fancypagestyle{primera}{\fancyhf{}\renewcommand{\headrulewidth}{0pt}
  \fancyfoot[L]{\raisebox{-0.32\height}{\includegraphics[height=0.78cm]{logo-gva.png}}\hspace{5mm}%
    \raisebox{-0.32\height}{\includegraphics[height=0.78cm]{logo-vermellar.png}}}
  \fancyfoot[R]{\scriptsize\color{muted}\href{https://creativecommons.org/licenses/by-nc-sa/4.0/deed.ca}{\ccbyncsa}\ \ CC BY-NC-SA 4.0 · fentmates.org\qquad\small\thepage}}
% la resta: el símbol a l'esquerra i el curs a la dreta, amb el títol al centre com abans
\fancyhead[L]{\raisebox{-0.3\height}{\includegraphics[height=0.62cm]{\fm@marca}}}
\fancyhead[R]{\small\bfseries\color{cursfosc}\fm@cursnom}
\renewcommand{\headrule}{\color{curs!45}\hrule height \headrulewidth}
\makeatother

%   \fmcurs{1eso}
% i, al començament del document:
%   \capcalera{Matemàtiques}{1r ESO}{Fitxa d'exercicis · Unitat 2}{Divisibilitat}
% Opcional abans de \input: \def\fmneed{14} (espai mínim abans de cada secció).
% =====================================================================
\makeatletter
\def\fm@logodir{../../logo/pdf}
\providecommand{\fmneed}{10}
\def\fm@cursnom{}

% clau del curs, color i color del text damunt del color
\newcommand{\fm@curs}[3]{\@namedef{fm@color@#1}{#2}\@namedef{fm@text@#1}{#3}}
\fm@curs{1eso}{E6B84A}{navy}
\fm@curs{2eso}{3B82F6}{white}
\fm@curs{3eso}{FF7A50}{white}
\fm@curs{4eso-a}{A78BFA}{white}
\fm@curs{4eso-b}{6D28D9}{white}
\fm@curs{1batx-mi}{9B1C1C}{white}
\fm@curs{1batx-ccss}{EF6B6B}{white}
\fm@curs{2batx-mii}{1E6B47}{white}
\fm@curs{2batx-ccss}{2BB3A3}{white}

\newcommand{\fmcurs}[1]{%
  \@ifundefined{fm@color@#1}{\PackageError{fentmates}{Curs desconegut: #1}{}}{}%
  \edef\fm@tmp{\noexpand\definecolor{curs}{HTML}{\@nameuse{fm@color@#1}}%
               \noexpand\colorlet{curstext}{\@nameuse{fm@text@#1}}}\fm@tmp
  \colorlet{cursfosc}{curs!72!black}%
  \def\fm@logo{\fm@logodir/vertical-#1.pdf}%
  \def\fm@marca{\fm@logodir/marca-#1.pdf}}

% ---------- capçalera de la primera pàgina ----------
% #1 matèria (petit), #2 curs, #3 tipus de document i unitat, #4 títol
\newcommand{\capcalera}[4]{%
  \gdef\fm@cursnom{#2}%
  \par\noindent
  \begin{tcolorbox}[colback=curs,colframe=curs,boxrule=0pt,arc=2.2mm,width=3.5cm,height=2.15cm,
      valign=center,left=3mm,right=2mm,top=1mm,bottom=1mm,nobeforeafter,box align=center]
    \color{curstext}{\footnotesize #1}\par\vspace{1pt}{\fontsize{19}{21}\selectfont\bfseries #2\par}%
  \end{tcolorbox}\hspace{5mm}%
  \begin{minipage}[c]{\dimexpr\linewidth-3.5cm-5mm-2.9cm\relax}
    {\color{cursfosc}\small\bfseries\MakeUppercase{#3}}\par\vspace{2pt}
    {\color{navy}\fontsize{20}{23}\selectfont\bfseries #4\par}
  \end{minipage}\hfill
  \begin{minipage}[c]{2.6cm}\raggedleft\includegraphics[width=2.6cm]{\fm@logo}\end{minipage}\par
  \vspace{2pt}}

% ---------- barres de secció del color del curs ----------
\renewcommand{\seccio}[2]{\par\needspace{\fmneed\baselineskip}\bigskip
  \begin{tcolorbox}[colback=curs,colframe=curs,coltext=curstext,arc=2mm,boxrule=0pt,left=3mm,right=3mm,top=1.2mm,bottom=1.2mm]
  {\large\bfseries #1\quad #2}
  \end{tcolorbox}\vspace{-2pt}}

% ---------- estrelles de dificultat del color del curs (a Batxillerat no se'n posen) ----------
\@ifundefined{estrelles}{}{%
  \renewcommand{\estrelles}[1]{{\color{curs}\ifcase#1\or\faStar\faStar[regular]\faStar[regular]\or\faStar\faStar\faStar[regular]\or\faStar\faStar\faStar\fi}}}

% ---------- capçalera i peu de les pàgines ----------
% primera pàgina: logos institucionals al peu
\fancypagestyle{primera}{\fancyhf{}\renewcommand{\headrulewidth}{0pt}
  \fancyfoot[L]{\raisebox{-0.32\height}{\includegraphics[height=0.78cm]{logo-gva.png}}\hspace{5mm}%
    \raisebox{-0.32\height}{\includegraphics[height=0.78cm]{logo-vermellar.png}}}
  \fancyfoot[R]{\scriptsize\color{muted}\href{https://creativecommons.org/licenses/by-nc-sa/4.0/deed.ca}{\ccbyncsa}\ \ CC BY-NC-SA 4.0 · fentmates.org\qquad\small\thepage}}
% la resta: el símbol a l'esquerra i el curs a la dreta, amb el títol al centre com abans
\fancyhead[L]{\raisebox{-0.3\height}{\includegraphics[height=0.62cm]{\fm@marca}}}
\fancyhead[R]{\small\bfseries\color{cursfosc}\fm@cursnom}
\renewcommand{\headrule}{\color{curs!45}\hrule height \headrulewidth}
\makeatother

%   \fmcurs{1eso}
% i, al començament del document:
%   \capcalera{Matemàtiques}{1r ESO}{Fitxa d'exercicis · Unitat 2}{Divisibilitat}
% Opcional abans de \input: \def\fmneed{14} (espai mínim abans de cada secció).
% =====================================================================
\makeatletter
\def\fm@logodir{../../logo/pdf}
\providecommand{\fmneed}{10}
\def\fm@cursnom{}

% clau del curs, color i color del text damunt del color
\newcommand{\fm@curs}[3]{\@namedef{fm@color@#1}{#2}\@namedef{fm@text@#1}{#3}}
\fm@curs{1eso}{E6B84A}{navy}
\fm@curs{2eso}{3B82F6}{white}
\fm@curs{3eso}{FF7A50}{white}
\fm@curs{4eso-a}{A78BFA}{white}
\fm@curs{4eso-b}{6D28D9}{white}
\fm@curs{1batx-mi}{9B1C1C}{white}
\fm@curs{1batx-ccss}{EF6B6B}{white}
\fm@curs{2batx-mii}{1E6B47}{white}
\fm@curs{2batx-ccss}{2BB3A3}{white}

\newcommand{\fmcurs}[1]{%
  \@ifundefined{fm@color@#1}{\PackageError{fentmates}{Curs desconegut: #1}{}}{}%
  \edef\fm@tmp{\noexpand\definecolor{curs}{HTML}{\@nameuse{fm@color@#1}}%
               \noexpand\colorlet{curstext}{\@nameuse{fm@text@#1}}}\fm@tmp
  \colorlet{cursfosc}{curs!72!black}%
  \def\fm@logo{\fm@logodir/vertical-#1.pdf}%
  \def\fm@marca{\fm@logodir/marca-#1.pdf}}

% ---------- capçalera de la primera pàgina ----------
% #1 matèria (petit), #2 curs, #3 tipus de document i unitat, #4 títol
\newcommand{\capcalera}[4]{%
  \gdef\fm@cursnom{#2}%
  \par\noindent
  \begin{tcolorbox}[colback=curs,colframe=curs,boxrule=0pt,arc=2.2mm,width=3.5cm,height=2.15cm,
      valign=center,left=3mm,right=2mm,top=1mm,bottom=1mm,nobeforeafter,box align=center]
    \color{curstext}{\footnotesize #1}\par\vspace{1pt}{\fontsize{19}{21}\selectfont\bfseries #2\par}%
  \end{tcolorbox}\hspace{5mm}%
  \begin{minipage}[c]{\dimexpr\linewidth-3.5cm-5mm-2.9cm\relax}
    {\color{cursfosc}\small\bfseries\MakeUppercase{#3}}\par\vspace{2pt}
    {\color{navy}\fontsize{20}{23}\selectfont\bfseries #4\par}
  \end{minipage}\hfill
  \begin{minipage}[c]{2.6cm}\raggedleft\includegraphics[width=2.6cm]{\fm@logo}\end{minipage}\par
  \vspace{2pt}}

% ---------- barres de secció del color del curs ----------
\renewcommand{\seccio}[2]{\par\needspace{\fmneed\baselineskip}\bigskip
  \begin{tcolorbox}[colback=curs,colframe=curs,coltext=curstext,arc=2mm,boxrule=0pt,left=3mm,right=3mm,top=1.2mm,bottom=1.2mm]
  {\large\bfseries #1\quad #2}
  \end{tcolorbox}\vspace{-2pt}}

% ---------- estrelles de dificultat del color del curs (a Batxillerat no se'n posen) ----------
\@ifundefined{estrelles}{}{%
  \renewcommand{\estrelles}[1]{{\color{curs}\ifcase#1\or\faStar\faStar[regular]\faStar[regular]\or\faStar\faStar\faStar[regular]\or\faStar\faStar\faStar\fi}}}

% ---------- capçalera i peu de les pàgines ----------
% primera pàgina: logos institucionals al peu
\fancypagestyle{primera}{\fancyhf{}\renewcommand{\headrulewidth}{0pt}
  \fancyfoot[L]{\raisebox{-0.32\height}{\includegraphics[height=0.78cm]{logo-gva.png}}\hspace{5mm}%
    \raisebox{-0.32\height}{\includegraphics[height=0.78cm]{logo-vermellar.png}}}
  \fancyfoot[R]{\scriptsize\color{muted}\href{https://creativecommons.org/licenses/by-nc-sa/4.0/deed.ca}{\ccbyncsa}\ \ CC BY-NC-SA 4.0 · fentmates.org\qquad\small\thepage}}
% la resta: el símbol a l'esquerra i el curs a la dreta, amb el títol al centre com abans
\fancyhead[L]{\raisebox{-0.3\height}{\includegraphics[height=0.62cm]{\fm@marca}}}
\fancyhead[R]{\small\bfseries\color{cursfosc}\fm@cursnom}
\renewcommand{\headrule}{\color{curs!45}\hrule height \headrulewidth}
\makeatother

\fmcurs{2batx-ccss}

\begin{document}
\thispagestyle{primera}

\capcalera{Matemáticas CCSS}{2\textsuperscript{o} Bach.}{Ficha de ejercicios · Unidad 5}{Cálculo de derivadas. Derivabilidad de una función}
\vspace{10pt}
Nombre y apellidos: \rule{6.6cm}{0.4pt}\hfill Grupo: \rule{1.6cm}{0.4pt}\hfill Fecha: \rule{2.2cm}{0.4pt}
\vspace{8pt}

\begin{tcolorbox}[colback=soft,colframe=soft,arc=2mm,left=3mm,right=3mm,top=1.5mm,bottom=1.5mm]
\small Los ejercicios marcados con {\color{teal}\bfseries\faUniversity\ PAU} son de las pruebas de acceso a la universidad de otras comunidades autónomas. Resuélvelos en el cuaderno mostrando todos los pasos y simplifica las derivadas.
\end{tcolorbox}

% =====================================================
\seccio{1}{Tasa de variación media e interpretación de la derivada}

\ex{1} Calcula la tasa de variación media de $f(x)=x^2-4x+1$ en los intervalos siguientes:\begin{tasks}(3)
  \task $[0,2]$
  \task $[1,4]$
  \task $[2,5]$
\end{tasks}\recorda{$\text{TVM}[a,b]=\dfrac{f(b)-f(a)}{b-a}$: es la pendiente de la recta secante que pasa por $(a,f(a))$ y $(b,f(b))$.}

\ex{1} Calcula la TVM de $g(x)=\sqrt x$ en $[1,4]$ y en $[4,9]$. ¿En cuál de los dos intervalos crece más rápidamente la función?

\ex{2} \label{ex:cost}El coste, en euros, de fabricar $x$ unidades de un producto es $C(x)=0{,}01x^2+5x+200$. Calcula la TVM del coste entre 100 y 200 unidades y entre 200 y 300 unidades. ¿Qué significan estos valores?

\ex{1} Responde a las preguntas siguientes:\begin{tasks}(1)
  \task Sabemos que $f(3)=2$ y $f'(3)=-4$. Escribe la ecuación de la recta tangente a la gráfica de $f$ en $x=3$.
  \task La recta tangente a $g$ en $x=1$ es $y=-2x+7$. ¿Cuánto valen $g(1)$ y $g'(1)$?
  \task Si $f'(a)=0$, ¿cómo es la recta tangente en $x=a$?
\end{tasks}\recorda{$f'(a)$ es la pendiente de la recta tangente en $x=a$. Recta tangente: $y-f(a)=f'(a)\,(x-a)$.}


% =====================================================
\seccio{2}{Derivada de una función por definición}

\ex{2} Calcula, aplicando la definición de derivada:\begin{tasks}(2)
  \task $f'(1)$ si $f(x)=3x^2-x$
  \task $f'(-1)$ si $f(x)=x^3-2x$
  \task $f'(3)$ si $f(x)=\dfrac{2}{x-1}$
  \task $f'(4)$ si $f(x)=\sqrt{2x+1}$
\end{tasks}\recorda{$f'(a)=\displaystyle\lim_{h\to0}\frac{f(a+h)-f(a)}{h}$. Simplifica $f(a+h)-f(a)$ hasta que aparezca el factor $h$; con raíces, multiplica por el conjugado.}

\ex{2} \pau{Navarra, EvAU extraordinaria 2023} Considera la función $g(x)=-x^2+4x+3$. Calcula la derivada de $g(x)$ en el punto $x=1$ aplicando la definición de derivada.

\ex{2} Calcula la función derivada aplicando la definición:\begin{tasks}(3)
  \task $f(x)=x^2-3x+2$
  \task $f(x)=\dfrac{1}{x+2}$
  \task $f(x)=\sqrt x$
\end{tasks}

\ex{2} Calcula por definición la derivada de $f(x)=x^2$ en $x=3$ y escribe la ecuación de la recta tangente a la gráfica en ese punto.


% =====================================================
\seccio{3}{Propiedades de la derivada: suma, producto y cociente}

\ex{1} Deriva estas funciones polinómicas:\begin{tasks}(2)
  \task $f(x)=5x^4-3x^3+2x-8$
  \task $f(x)=\dfrac{x^5}{5}-\dfrac{2x^3}{3}+x$
  \task $f(x)=\dfrac{3x^2-2x+1}{4}$
  \task $f(x)=-x^6+\dfrac x2-7$
\end{tasks}\recorda{$(k)'=0$, \ $(x^n)'=n\,x^{n-1}$, \ $(f\pm g)'=f'\pm g'$, \ $(k\cdot f)'=k\cdot f'$.}

\ex{1} Deriva aplicando la regla del producto y simplifica:\begin{tasks}(3)
  \task $f(x)=(2x-1)(x^2+3)$
  \task $f(x)=(x^3+x)(1-x^2)$
  \task $f(x)=x^2(x-1)(x+1)$
\end{tasks}\recorda{$(f\cdot g)'=f'\cdot g+f\cdot g'$. \ Cuidado: $(f\cdot g)'\neq f'\cdot g'$.}

\ex{2} Deriva aplicando la regla del cociente y simplifica:\begin{tasks}(2)
  \task $f(x)=\dfrac{x+3}{x-2}$
  \task $f(x)=\dfrac{2x^2}{x+1}$
  \task $f(x)=\dfrac{x^2-1}{x^2+1}$
  \task $f(x)=\dfrac{3}{x^2-4}$
\end{tasks}\recorda{$\left(\dfrac fg\right)'=\dfrac{f'\cdot g-f\cdot g'}{g^2}$. Simplifica el numerador, pero deja el denominador sin desarrollar.}

\ex{1} Dada $f(x)=x^4-2x^3+5x$, calcula $f'(x)$, $f''(x)$ y $f''(1)$.

\ex{2} El coste de fabricar $x$ unidades es $C(x)=0{,}01x^2+5x+200$ euros (ejercicio \ref{ex:cost}). Calcula el coste marginal $C'(x)$ y $C'(100)$. ¿Qué significa? Compáralo con la TVM del ejercicio \ref{ex:cost}.


% =====================================================
\seccio{4}{Derivadas de las funciones elementales}

\ex{1} Escribe como potencia y deriva:\begin{tasks}(3)
  \task $f(x)=\sqrt{x^3}$
  \task $f(x)=\dfrac{1}{x^2}$
  \task $f(x)=\sqrt[3]{x}$
  \task $f(x)=\dfrac{6}{\sqrt x}$
  \task $f(x)=x^2\sqrt x$
  \task $f(x)=\dfrac{2}{3x^3}$
\end{tasks}\recorda{$\sqrt[n]{x^m}=x^{m/n}$ y $\dfrac{1}{x^n}=x^{-n}$. La regla $(x^n)'=n\,x^{n-1}$ vale para cualquier exponente.}

\ex{1} Deriva estas funciones exponenciales y logarítmicas:\begin{tasks}(3)
  \task $f(x)=2e^x-3\ln x$
  \task $f(x)=4^x$
  \task $f(x)=\log_3x$
  \task $f(x)=x\,e^x$
  \task $f(x)=\dfrac{e^x}{x}$
  \task $f(x)=x^2\ln x$
\end{tasks}\recorda{$(e^x)'=e^x$, \ $(a^x)'=a^x\ln a$, \ $(\ln x)'=\dfrac1x$, \ $(\log_ax)'=\dfrac{1}{x\ln a}$.}

\ex{2} Deriva estas funciones trigonométricas:\begin{tasks}(3)
  \task $f(x)=3\sin x-2\cos x$
  \task $f(x)=x\sin x$
  \task $f(x)=\tan x-x$
  \task $f(x)=\sin x\cos x$
  \task $f(x)=\arctan x-x$
  \task $f(x)=x\arcsin x$
\end{tasks}\recorda{$(\sin x)'=\cos x$, \ $(\cos x)'=-\sin x$, \ $(\tan x)'=1+\tan^2x$, \ $(\arcsin x)'=\frac{1}{\sqrt{1-x^2}}$, \ $(\arctan x)'=\frac{1}{1+x^2}$.}


% =====================================================
\seccio{5}{Regla de la cadena}

\ex{2} Deriva estas potencias y raíces de funciones:\begin{tasks}(2)
  \task $f(x)=(3x^2-1)^4$
  \task $f(x)=\sqrt{x^2-2x}$
  \task $f(x)=\dfrac{1}{(x^2+1)^3}$
  \task $f(x)=\sqrt[3]{2x+5}$
\end{tasks}\recorda{Regla de la cadena: $[f(g(x))]'=f'(g(x))\cdot g'(x)$. Por ejemplo, $[g^n]'=n\,g^{n-1}\cdot g'$ y $\left[\sqrt g\right]'=\dfrac{g'}{2\sqrt g}$.}

\ex{2} Deriva estas funciones compuestas exponenciales y logarítmicas:\begin{tasks}(3)
  \task $f(x)=e^{3x-1}$
  \task $f(x)=2^{x^2}$
  \task $f(x)=\ln(x^2-3x)$
  \task $f(x)=\log(5x)$
  \task $f(x)=\ln\sqrt x$
  \task $f(x)=e^{\sqrt x}$
\end{tasks}\recorda{$\left[e^{g}\right]'=e^{g}\cdot g'$, \ $\left[a^{g}\right]'=a^{g}\cdot g'\cdot\ln a$, \ $\left[\ln g\right]'=\dfrac{g'}{g}$. Antes de derivar un logaritmo, aplica sus propiedades si simplifican.}

\ex{2} Deriva estas funciones trigonométricas compuestas:\begin{tasks}(3)
  \task $f(x)=\sin(5x+1)$
  \task $f(x)=\cos(x^2)$
  \task $f(x)=\tan(3x)$
  \task $f(x)=\sin^3x$
  \task $f(x)=\arctan(2x)$
  \task $f(x)=\arcsin(x^2)$
\end{tasks}

\ex{2} Deriva y simplifica (hay más de una regla en cada función):\begin{tasks}(2)
  \task $f(x)=x\,e^{-x^2}$
  \task $f(x)=\ln\dfrac{x+1}{x^2+1}$
  \task $f(x)=(x-1)^2(x+2)^3$
  \task $f(x)=\dfrac{e^x}{e^x+1}$
\end{tasks}

\ex{2} \pau{Murcia, EBAU septiembre 2020} Calcula las derivadas de las funciones siguientes:\begin{tasks}(2)
  \task $f(x)=(x^2-2)\ln x$
  \task $f(x)=e^{4x^3+2}$
\end{tasks}

\ex{2} \pau{Murcia, EBAU junio 2018} Halla las derivadas de las funciones siguientes:\begin{tasks}(2)
  \task $g(x)=x^2\ln(x^2)$
  \task $h(x)=e^{-3x+x^2}$
\end{tasks}

\ex{2} \pau{Andalucía, PEvAU ordinaria 2021} Calcula la derivada de las funciones siguientes:\begin{tasks}(2)
  \task $f(x)=\ln\left(\dfrac{x-1}{x+1}\right)$
  \task $g(x)=x^3\cdot e^{2x^2}$
\end{tasks}

\ex{2} \pau{Andalucía, ordinaria (suplente) 2025} Calcula la derivada de las funciones siguientes:\begin{tasks}(2)
  \task $g(x)=\ln\left(\dfrac{x^2-1}{x^2+1}\right)$
  \task $h(x)=(2x-1)\,e^{x^2-x}$
\end{tasks}

\ex{3} \pau{Andalucía, PEvAU ordinaria 2024} Calcula la derivada de las funciones siguientes:\begin{tasks}(2)
  \task $f(x)=(x^2+2)^3\cdot e^{-2x}$
  \task $g(x)=\dfrac{\ln(1-x^3)}{(1-2x^2)^2}$
\end{tasks}

\ex{2} \pau{Murcia, EBAU septiembre 2018} Dada la función $f(x)=5x^3e^{2x}+\dfrac{1}{x^2+1}$:\begin{tasks}(1)
  \task Calcula $f'(0)$.
  \task Halla la ecuación de la recta tangente a la gráfica de la función en el punto $(0,1)$.
\end{tasks}

\ex{2} Halla la ecuación de la recta tangente a la gráfica de cada función en el punto indicado:\begin{tasks}(1)
  \task \pau{Murcia, EBAU julio 2020} $f(x)=\dfrac{1}{x+1}$ en $x=1$.
  \task \pau{Navarra, EvAU ordinaria 2023} $f(x)=\sqrt{5+x^2}$ en $x=-2$.
  \task \pau{Navarra, EvAU ordinaria 2022} $f(x)=\dfrac{x}{x^2-4}$ en $x=1$.
\end{tasks}

\ex{3} \pau{Islas Baleares, PBAU extraordinaria 2021} Considera la función $f(x)=\dfrac3x+8$.\begin{tasks}(1)
  \task Halla los puntos de la gráfica en los que la recta tangente es paralela a la recta $3x+4y+5=0$.
  \task Calcula las ecuaciones de las rectas tangentes en los puntos hallados en el apartado anterior.
\end{tasks}

\ex{2} Halla los parámetros en cada caso:\begin{tasks}(1)
  \task \pau{Madrid, EvAU extraordinaria 2022} $a$ y $b$ para que $f(x)=ax+\dfrac bx$ verifique $f(2)=4$ y $f'(2)=0$.
  \task \pau{Murcia, EBAU septiembre 2017} $a$ y $b$ para que $f(x)=\dfrac{ax+b}{x^2+1}$ cumpla $f(0)=1$ y $f'(0)=1$.
  \task \pau{Murcia, EBAU septiembre 2018} $a$ y $b$ para que $f(x)=x^3\ln(2x+5)+ax+b$ cumpla $f(0)=2$ y $f'(0)=1$.
  \task \pau{Andalucía, PEvAU ordinaria 2024} $a$ y $b$ para que sea horizontal la recta tangente a la gráfica de $h(x)=x^3+ax^2+3x+b$ en el punto $P(1,2)$.
\end{tasks}


% =====================================================
\seccio{6}{Continuidad y derivabilidad. Derivadas laterales}

\ex{2} Calcula las derivadas laterales en el punto indicado y di si la función es derivable en él:\begin{tasks}(2)
  \task $f(x)=|x+1|$ en $x=-1$
  \task $f(x)=|x^2-9|$ en $x=3$
  \task $f(x)=x\,|x-1|$ en $x=1$
  \task $f(x)=\sqrt[3]{x^2}$ en $x=0$
\end{tasks}\recorda{$f'(a^-)=\displaystyle\lim_{h\to0^-}\frac{f(a+h)-f(a)}{h}$ y $f'(a^+)=\displaystyle\lim_{h\to0^+}\frac{f(a+h)-f(a)}{h}$. \ $f$ es derivable en $a$ si las dos existen, son finitas y coinciden. Un valor absoluto se escribe primero a trozos.}

\ex{1} Razona si son verdaderas o falsas las afirmaciones siguientes. Si son falsas, pon un contraejemplo:\begin{tasks}(1)
  \task Si $f$ es continua en $x=a$, entonces es derivable en $x=a$.
  \task Si $f$ no es continua en $x=a$, entonces no es derivable en $x=a$.
  \task Si $f$ no es derivable en $x=a$, entonces no es continua en $x=a$.
  \task Si en una función a trozos las derivadas de los dos trozos valen lo mismo en el punto de cambio $a$, la función es derivable en $a$.
\end{tasks}\recorda{Derivable $\Rightarrow$ continua. El recíproco no es cierto.}

\ex{2} Observa la gráfica de la función $f$ e indica en qué puntos no es continua y en cuáles no es derivable. Justifícalo. (En $x=0$ la recta y la curva se unen sin formar pico.)

\begin{center}\begin{tikzpicture}[scale=0.62,baseline=(current bounding box.center)]\begin{scope}\draw[muted!30,very thin] (-4,-1) grid (5,4);\draw[-latex,thick] (-4,0)--(5.4,0);\draw[-latex,thick] (0,-1)--(0,4.4);\clip (-4,-1) rectangle (5,4);\draw[accent,line width=1.5pt] (-4,1)--(-2,3)--(0,1);\draw[accent,line width=1.5pt,smooth,samples=60,domain=0:2] plot (\x,{1-\x+\x*\x/2});\draw[accent,line width=1.5pt] (2,3)--(5,3);\filldraw[fill=white,draw=accent,line width=1.1pt] (2,1) circle (0.13);\fill[accent] (2,3) circle (0.13);\node[below,font=\scriptsize] at (-2,0) {$-2$};\node[below,font=\scriptsize] at (2,0) {$2$};\node[below,font=\scriptsize] at (1,0) {$1$};\node[left,font=\scriptsize] at (0,1) {$1$};\node[left,font=\scriptsize] at (0,3) {$3$};\end{scope}\end{tikzpicture}\end{center}


% =====================================================
\seccio{7}{Derivabilidad de funciones a trozos}

\ex{2} Estudia la continuidad y la derivabilidad de las funciones siguientes:\begin{tasks}(2)
  \task $f(x)=\begin{cases}e^x&\text{si }x<0\\[3pt]x+1&\text{si }x\geq0\end{cases}$
  \task $g(x)=\begin{cases}3-x^2&\text{si }x<1\\[3pt]\dfrac2x&\text{si }x\geq1\end{cases}$
\end{tasks}\recorda{1) Continuidad en el punto de cambio (si no es continua, no es derivable). 2) Derivamos cada trozo \textbf{sin el igual}. 3) Comparamos $f'(a^-)$ y $f'(a^+)$.}

\ex{2} \pau{Andalucía, PEvAU extraordinaria 2024} Considera la función\[f(x)=\begin{cases}-x^2+4x+3&\text{si }x<4\\[3pt]2x-5&\text{si }x\geq4\end{cases}\]Estudia su continuidad y su derivabilidad.

\ex{2} \pau{Andalucía, ordinaria (suplente) 2025} Considera la función\[f(x)=\begin{cases}10+\dfrac{5x}{2}&\text{si }x\leq-2\\[3pt]x^2+1&\text{si }-2<x<2\\[3pt]10-\dfrac{5x}{2}&\text{si }x\geq2\end{cases}\]Estudia la continuidad y la derivabilidad de $f$ en el punto de abscisa $x=-2$.

\ex{2} \pau{Andalucía, PEvAU ordinaria 2024} La velocidad del viento, en km/h, prevista para cierto día en una zona de Sierra Nevada viene dada por la función\[v(t)=\begin{cases}t^2-8t+60&\text{si }0\leq t\leq10\\[3pt]-t^2+32t-140&\text{si }10<t\leq24\end{cases}\]donde $t$ es el tiempo en horas. Comprueba que la función $v$ es continua y derivable.

\ex{2} \pau{Navarra, EvAU ordinaria 2024} Sea la función\[f(x)=\begin{cases}-x^2-4x&\text{si }x<-1\\[3pt]4-x&\text{si }-1\leq x<1\\[3pt]x^2-6x+8&\text{si }x\geq1\end{cases}\]\begin{tasks}(1)
  \task Estudia la continuidad de $f(x)$, clasificando sus puntos de discontinuidad.
  \task Estudia la derivabilidad de $f(x)$.
\end{tasks}

\ex{2} \pau{Andalucía, PEvAU extraordinaria 2021} Considera la función\[f(x)=\begin{cases}2^{x+1}&\text{si }x<0\\[3pt]x^2-2x&\text{si }x\geq0\end{cases}\]Estudia la continuidad y la derivabilidad de la función $f$ en su dominio.

\ex{3} \pau{Andalucía, PEvAU extraordinaria 2021} El número de diagnosticados por PCR en Andalucía, en miles de personas, se aproxima por la función\[f(t)=\begin{cases}-t^2+2t-0{,}3&\text{si }0{,}2\leq t\leq1{,}8\\[3pt]0{,}1t-0{,}12&\text{si }1{,}8<t\leq5\\[3pt]-0{,}5t^2+8{,}3t-28{,}62&\text{si }5<t\leq10\end{cases}\]donde $t$ es el tiempo en meses. Estudia la continuidad y la derivabilidad de la función $f$ en su dominio.

\ex{3} Escribe a trozos y estudia la derivabilidad de $f(x)=|x-1|+|x+1|$ y de $g(x)=|x^2-4x|$.


% =====================================================
\seccio{8}{Funciones a trozos con parámetros}

\ex{2} \pau{Andalucía, PEvAU ordinaria 2022} Considera la función\[f(x)=\begin{cases}6x-3&\text{si }x\leq1\\[3pt]ax^2+bx+2&\text{si }x>1\end{cases}\]con $a$ y $b$ números reales. Determina los valores de $a$ y $b$ para que $f$ sea continua y derivable en todo su dominio.\recorda{En el punto de cambio planteamos dos ecuaciones: continuidad ($\lim_{x\to a^-}f=\lim_{x\to a^+}f=f(a)$) y derivabilidad ($f'(a^-)=f'(a^+)$).}

\ex{2} \pau{Andalucía, PEvAU extraordinaria 2020} Considera la función\[f(x)=\begin{cases}2+\dfrac{a}{x-1}&\text{si }x<0\\[3pt]a+be^x&\text{si }x\geq0\end{cases}\]Calcula los valores de $a$ y $b$ para que la función sea continua y derivable en su dominio.

\ex{2} \pau{Islas Baleares, PBAU ordinaria 2022} Consideremos la función a trozos\[f(x)=\begin{cases}x^2-3x+2&\text{si }x<0\\[3pt]x^2+ax+2&\text{si }x\geq0\end{cases}\]Calcula los valores de $a$ para que $f(x)$ sea continua y derivable.

\ex{2} \pau{Islas Baleares, PBAU ordinaria 2021} Considera la función a trozos\[f(x)=\begin{cases}x^2-3x+2&\text{si }x<0\\[3pt]e^{ax}+1&\text{si }x\geq0\end{cases}\]Calcula los valores de $a$ para que $f$ sea continua y derivable.

\ex{2} \pau{Islas Baleares, PBAU extraordinaria 2021} Considera la función definida a trozos\[f(x)=\begin{cases}-4x+a&\text{si }x\leq-2\\[3pt]x^2-5&\text{si }-2<x<1\\[3pt]-7x+3&\text{si }x\geq1\end{cases}\]\begin{tasks}(1)
  \task Calcula los valores de $a$ para que $f(x)$ sea continua.
  \task ¿Es $f(x)$ derivable para $a=1$?
\end{tasks}

\ex{3} \pau{Aragón, EvAU ordinaria 2022} Siendo $a$, $b$ parámetros reales, se considera la función\[f(x)=\begin{cases}x^2+2\sqrt3&\text{si }x\leq0\\[3pt]\sqrt{ax+b}&\text{si }0<x\leq3\\[3pt]\dfrac72-\dfrac x6&\text{si }x>3\end{cases}\]\begin{tasks}(1)
  \task Determina el valor de los parámetros para que $f(x)$ sea continua.
  \task Para estos valores, analiza si $f(x)$ es derivable en $x=0$ y en $x=3$.
\end{tasks}

\ex{2} \pau{Andalucía, ordinaria (suplente) 2025} El nivel de concentración de un alumno durante un examen de 5 horas viene dado por la función\[f(t)=\begin{cases}-t^2+2t+10&\text{si }0\leq t\leq2{,}5\\[3pt]t^2+at+b&\text{si }2{,}5<t\leq5\end{cases}\]donde $t$ es el tiempo en horas y $a$ y $b$ son números reales. ¿Con qué nivel de concentración empieza el alumno el examen? Determina los valores de $a$ y $b$ para que la función sea continua y derivable en $t=2{,}5$.

\ex{2} Halla los parámetros para que cada función sea derivable en $\mathbb R$:\begin{tasks}(2)
  \task $f(x)=\begin{cases}ax^2+bx+1&\text{si }x\leq2\\[3pt]3x-3&\text{si }x>2\end{cases}$
  \task $g(x)=\begin{cases}e^{2x}&\text{si }x\leq0\\[3pt]x^2+ax+b&\text{si }x>0\end{cases}$
\end{tasks}

\vfill
\begin{center}\color{muted}\small\faCheckCircle[regular]\ \ ¡Ánimo! Comprueba tus resultados con la calculadora de derivadas de los apuntes interactivos.\end{center}

\end{document}
